\subsection{Voting Rights Act Requirements}

Section 2 of the Voting Rights Act, as amended in 1982, prohibits any voting practice that ``results in a denial or abridgement of the right to vote on account of race or color'' \citep{vra1982}. Courts apply the \emph{Gingles} test \citep{thornburg1986} to determine whether redistricting plans violate Section 2:

\begin{enumerate}
\item The minority group must be sufficiently large and geographically compact to constitute a majority in a single-member district.
\item The minority group must be politically cohesive.
\item The majority votes sufficiently as a bloc to enable it to defeat the minority's preferred candidate.
\end{enumerate}

When these conditions are met, jurisdictions must create majority-minority (MM) districts where the minority voting-age population (VAP) exceeds 50\%. In practice, courts often require higher thresholds (55-65\%) to account for lower registration and turnout rates \citep{katz2020}.

\subsection{Algorithmic Redistricting}

Recent work on algorithmic redistricting emphasizes compactness, contiguity, and population balance as core principles \citep{duchin2018, fifield2020, deford2021}. Common approaches include:

\textbf{Recursive Bisection}: Hierarchically partitioning geography into districts by repeatedly splitting regions in two \citep{barnes1982, karypis1998}. METIS graph partitioning \citep{karypis1998metis} minimizes edge cuts (compactness proxy) while balancing population.

\textbf{Markov Chain Monte Carlo (MCMC)}: Generating ensembles of valid redistricting plans through random perturbations \citep{duchin2018, deford2021}. Enables statistical analysis of partisan bias and outlier detection.

\textbf{Optimization}: Direct optimization of redistricting criteria using integer programming \citep{hess1965, shirabe2005} or metaheuristics \citep{bacao2005}.

However, existing work focuses primarily on compactness and partisan fairness, with limited attention to VRA compliance as a hard constraint.

\subsection{Multi-Constraint Graph Partitioning}

METIS supports multi-constraint partitioning where vertices have multiple weights (e.g., population + minority VAP) and optimization balances all constraints simultaneously \citep{karypis1998metis}. Target partition weights (\texttt{tpwgts}) specify desired distribution for each constraint. Per-constraint imbalance tolerance (\texttt{ubvec}) allows relaxation of specific constraints.

Prior work uses multi-constraint partitioning for heterogeneous computing \citep{catalyurek2011} and scientific computing \citep{hendrickson2000}, but not for VRA-aware redistricting.

\subsection{VRA Compliance in Practice}

Real-world redistricting combines algorithmic tools with political negotiation and legal review \citep{pildes2002}. States with VRA requirements often create ``packed'' MM districts with 60-70\% minority VAP, sometimes at the expense of compactness \citep{lublin1997}.

Recent work \citep{duchin2020alabama, chen2020texas} analyzes VRA compliance in specific states but lacks systematic comparison of algorithmic approaches or identification of feasibility thresholds.

\subsection{Gap in Literature}

No prior work systematically evaluates whether principled multi-constraint optimization can achieve VRA compliance across multiple covered states. Our work fills this gap by:
(1) Testing four optimization approaches on five VRA states,
(2) Identifying critical demographic thresholds for feasibility,
(3) Quantifying the tradeoff between VRA compliance and compactness.
